\begingroup
\raggedright
\subsection{Formatted Output}
The table below is for \texttt{printf}. Use \texttt{\%.3e} for
three digits after the decimal point in scientific notation, or
\texttt{\%.2f} for two digits in fixed notation.
\texttt{\%e}/\texttt{\%E} expect \texttt{double};
\texttt{\%Le}/\texttt{\%LE} expect \texttt{long double}.
For \texttt{scanf}, use \texttt{\%f} for \texttt{float*},
\texttt{\%lf} for \texttt{double*}, and \texttt{\%Lf} for
\texttt{long double*}.

  \begin{tabular}{|c|p{6cm}|c|}
  \hline
  \textbf{specifier} & \textbf{Output} & \textbf{Example} \\ \hline

  d or i & Signed decimal integer & 392 \\ \hline
  u & Unsigned decimal integer & 7235 \\ \hline
  o & Unsigned octal & 610 \\ \hline
  x & Unsigned hexadecimal integer (lowercase) & 7fa \\ \hline
  X & Unsigned hexadecimal integer (uppercase) & 7FA \\ \hline
  f & Decimal floating point, lowercase & 392.65 \\ \hline
  F & Decimal floating point, uppercase & 392.65 \\ \hline
  e & Scientific notation (mantissa/exponent), lowercase & 3.9265e+2 \\ \hline
  E & Scientific notation (mantissa/exponent), uppercase & 3.9265E+2 \\ \hline
  g & General format (significant digits): \%e or \%f & 392.65 \\ \hline
  G & General format (significant digits): \%E or \%F & 392.65 \\ \hline
  a & Hexadecimal floating point, lowercase & -0xc.90fep-2 \\ \hline
  A & Hexadecimal floating point, uppercase & -0XC.90FEP-2 \\ \hline
  c & Character & a \\ \hline
  s & String of characters & sample \\ \hline
  p & Pointer address & b8000000 \\ \hline
  n & Nothing printed. The argument must be a pointer to a \texttt{signed int}. The number of characters written so far is stored. & \\ \hline
  \% & A \% followed by another \% prints a single \% & \% \\ \hline

  \end{tabular}
\endgroup
